\documentclass[11pt,a4paper]{article}

\usepackage[utf8]{inputenc}
\usepackage[T2A]{fontenc}
\usepackage[ukrainian,english]{babel}
\usepackage[top=22mm,bottom=22mm,left=22mm,right=22mm]{geometry}
\usepackage{hyperref}
\usepackage{fancyhdr}
\usepackage{parskip}
\usepackage{tcolorbox}
\usepackage{enumitem}
\usepackage{listings}
\usepackage{xcolor}

\tcbuselibrary{skins,breakable}

\hypersetup{
    pdftitle={Workshop 04 — Skills: Handout},
    pdfauthor={Vadym Bondarenko},
    colorlinks=true,
    linkcolor=blue,
    urlcolor=cyan,
}

\pagestyle{fancy}
\fancyhf{}
\fancyhead[L]{\small\textit{Workshop 04 — Skills}}
\fancyhead[R]{\small\thepage}
\fancyfoot[C]{\small\textit{vcdev.me · 2026}}
\renewcommand{\headrulewidth}{0.4pt}

\lstdefinelanguage{yaml}{
    keywords={true,false,null},
    sensitive=false,
    comment=[l]{\#},
    morestring=[b]',
    morestring=[b]"
}

\lstset{
    basicstyle=\ttfamily\small,
    breaklines=true,
    frame=single,
    backgroundcolor=\color{gray!10},
    rulecolor=\color{gray!40},
    xleftmargin=8pt,
    xrightmargin=8pt,
    aboveskip=6pt,
    belowskip=6pt,
}

\title{%
  {\Huge\bfseries Workshop 04: Skills}\\[0.6em]
  {\large Від ідеї до продакшну — handout}
}
\author{Vadym Bondarenko \\ \small\href{mailto:vadym.bondarenko@vcdev.me}{vadym.bondarenko@vcdev.me}}
\date{2026}

\begin{document}

\maketitle

\begin{abstract}
Пост-воркшоп референс для самостійного перечитування. Не повторює слайди — сухе зведення команд, frontmatter-полів, чек-листів і decision-trees. Усі факти посилаються на офіційні docs (\href{https://code.claude.com/docs/en/skills}{code.claude.com}).
\end{abstract}

\tableofcontents
\newpage

\section{Що таке skill}

Markdown-файл (\texttt{SKILL.md}) з YAML-frontmatter. Лежить у теці. Claude бачить frontmatter усіх skill-ів завжди; тіло вантажить лише при тригері.

\subsection{Локації}

\begin{tabular}{lll}
\hline
\textbf{Scope} & \textbf{Path} & \textbf{Назва інвокації} \\
\hline
Personal & \texttt{\textasciitilde/.claude/skills/<name>/SKILL.md} & \texttt{/<name>} \\
Project & \texttt{.claude/skills/<name>/SKILL.md} & \texttt{/<name>} \\
Plugin & \texttt{<plugin>/skills/<name>/SKILL.md} & \texttt{/<plugin>:<name>} \\
\hline
\end{tabular}

\textbf{Precedence:} Enterprise $>$ Personal $>$ Project $>$ Plugin.

\section{Frontmatter cheatsheet}

\subsection{Критичні поля}

\begin{lstlisting}[language=yaml]
---
name: skill-name              # kebab-case, ≤64, default = directory
description: Use when user asks ...   # головний trigger
allowed-tools: [Bash, Read]   # pre-approve без per-use prompts
---
\end{lstlisting}

\textbf{Бюджет:} \texttt{description} + \texttt{when\_to\_use} разом $\le$ 1536 символів.

\subsection{Контроль інвокації}

\begin{lstlisting}[language=yaml]
disable-model-invocation: true   # лише /name від юзера, Claude не може
user-invocable: false            # Claude може, юзер не бачить у /
paths: ["**/*.rs"]               # тригер тільки коли файли матчать
\end{lstlisting}

\subsection{Ізоляція}

\begin{lstlisting}[language=yaml]
context: fork                    # subagent, окремий контекст
agent: Explore                   # тип subagent (Explore | Plan | general-purpose | custom)
model: claude-sonnet             # override моделі поки skill активний
effort: high                     # low | medium | high | xhigh | max
\end{lstlisting}

\subsection{Аргументи}

\begin{lstlisting}[language=yaml]
argument-hint: [good-sha] [bad-sha]   # автокомпліт-підказка
arguments: [good, bad]                 # named positional, доступні як $good, $bad
shell: bash                            # bash | powershell — для !`команда`
\end{lstlisting}

\section{Anatomy теки}

\begin{lstlisting}
my-skill/
├── SKILL.md            # обов'язковий entry
├── scripts/            # опційно — bash/python utility
├── references/         # опційно — doc-и для on-demand завантаження
└── assets/             # опційно — шаблони, ікони
\end{lstlisting}

\textbf{Розмір:} цільтесь у $\le$500 рядків у \texttt{SKILL.md}. Більше --- виносьте у \texttt{references/}.

\section{Progressive disclosure}

3 рівні завантаження, кожен наступний дорожчий:

\begin{enumerate}
\item \textbf{Метадата} (name + description) — завжди в контексті
\item \textbf{Тіло SKILL.md} — при тригері
\item \textbf{Supporting files} — лише при покликанні з SKILL.md
\end{enumerate}

\textbf{Як розбивати:}

\begin{lstlisting}
cloud-deploy/
├── SKILL.md             # decision tree
└── references/
    ├── aws.md           # деталі AWS
    ├── gcp.md           # деталі GCP
    └── azure.md         # деталі Azure
\end{lstlisting}

У \texttt{SKILL.md} обов'язково \textbf{markdown-link} (не plain text):

\begin{lstlisting}[language=]
For AWS: read [aws.md](references/aws.md)
\end{lstlisting}

\section{Description: правила тригерів}

\subsection{Що робить хорошу description}

\begin{itemize}
\item Front-load trigger у перших 50 символах
\item Action verbs (\texttt{Use when}, \texttt{Diagnose}, \texttt{Refactor})
\item Перерахуй \textbf{синоніми} (як юзер може сказати по-різному)
\item Явно \textbf{коли} (контексти, тригерні фрази), не лише \textbf{що}
\item «Pushy» якщо skill часто недо-тригериться
\end{itemize}

\subsection{Red flags}

\begin{itemize}
\item Занадто загально: \emph{«Helpful tips»}
\item Без action verb: \emph{«Code explanations»}
\item Занадто вузько (overfit): \emph{«Fix bugs in TypeScript files only»}
\item Закопаний use case (description починається з \emph{«This skill provides»})
\item Дублює \texttt{name}
\end{itemize}

\section{Plugin packaging}

\subsection{Структура}

\begin{lstlisting}
my-plugin/
├── .claude-plugin/
│   └── plugin.json     # маніфест
└── skills/
    └── skill-name/
        └── SKILL.md
\end{lstlisting}

\textbf{Pitfall:} НЕ кладіть \texttt{skills/}, \texttt{commands/}, \texttt{hooks/} \textbf{всередину} \texttt{.claude-plugin/}.

\subsection{Мінімальний \texttt{plugin.json}}

\begin{lstlisting}[language=]
{
  "name": "my-plugin",
  "description": "Brief description",
  "version": "0.1.0",
  "author": { "name": "You" }
}
\end{lstlisting}

\textbf{Versioning:}
\begin{itemize}
\item Set explicit \texttt{version} → юзери оновлюються лише при bump
\item Skip → commit SHA = версія, кожен коміт оновлює юзерів
\end{itemize}

\textbf{Інвокація:} \texttt{/my-plugin:skill-name} (namespaced).

\section{Debug: skill не тригериться}

6-крокова діагностика:

\begin{enumerate}
\item \textbf{Skill в \texttt{/}-меню?} Ні → перевір шлях \texttt{\textasciitilde/.claude/skills/<name>/SKILL.md}
\item \textbf{Перефразуй точно під description.} Спрацювало → tuning description
\item \textbf{Грепни \texttt{disable-model-invocation} / \texttt{user-invocable}}
\item \textbf{Перевір \texttt{paths}} — може skill активується лише на матчинг файлах
\item \textbf{Manual \texttt{/skill-name} працює?} Так → skill loaded, проблема matching
\item \textbf{Запит надто простий?} Claude не тригерить skill для тривіальних задач (\textit{by design}).
\end{enumerate}

\subsection{Бюджет descriptions}

$\sim$1\% контекстного вікна = $\sim$8K символів за замовчуванням. 100+ skill-ів з довгими description --- обріжуть.

\textbf{Підняти:} \texttt{SLASH\_COMMAND\_TOOL\_CHAR\_BUDGET=16000}

\section{Production-чек-ліст}

Перед публікацією:

\begin{itemize}[leftmargin=2em]
\item[$\square$] \texttt{name} kebab-case, унікальний у твоєму неймспейсі
\item[$\square$] \texttt{description} front-load trigger, action verb, синоніми
\item[$\square$] \texttt{SKILL.md} $\le$ 500 рядків, важке у \texttt{references/}
\item[$\square$] Markdown-links на supporting файли (не plain text)
\item[$\square$] \texttt{allowed-tools} лише потрібне, не all
\item[$\square$] \texttt{disable-model-invocation: true} для всього з side-effects
\item[$\square$] Тестовано неявно (семантичні запити, не лише \texttt{/skill-name})
\item[$\square$] \texttt{README.md} у репо: як встановити, що робить
\item[$\square$] \texttt{version} у \texttt{plugin.json} якщо semver
\end{itemize}

\section{Команди для вправ}

\subsection{Вправа 1 --- встановити skill}

\begin{lstlisting}
mkdir -p ~/.claude/skills/git-status-summary
cp starter/SKILL.md ~/.claude/skills/git-status-summary/SKILL.md
# перезапусти Claude Code
\end{lstlisting}

\subsection{Вправа 4 --- встановити плагін}

\begin{lstlisting}
claude --plugin-dir ./my-git-toolkit
# далі:
/my-git-toolkit:git-status-summary
/my-git-toolkit:git-bisect-helper <good-sha> <bad-sha>
\end{lstlisting}

\section{Resources}

\begin{itemize}
\item \href{https://code.claude.com/docs/en/skills}{code.claude.com/docs/en/skills} --- Skills overview
\item \href{https://code.claude.com/docs/en/plugins}{code.claude.com/docs/en/plugins} --- Plugins
\item \href{https://agentskills.io}{agentskills.io} --- Open Agent Skills standard
\item Цей репо: \texttt{workshops/04-skills/exercises/} і \texttt{solutions/}
\end{itemize}

\end{document}
